\section{The PT Framework}
\label{sec:cadre_pt}

\subsection{The Fundamental Parameter}
\label{sec:parametre}

The Theory of Persistence (PT) takes as unique arithmetic axiom the
theorem~$\Tone$ \cite{PT_M1}:
\[
\sper \;=\; \tfrac{1}{2}.
\]
This parameter plays three distinct roles in the PT corpus, which will be
progressively identified in the course of the article:
\begin{itemize}[nosep,leftmargin=*]
\item \textbf{arithmetic}: $\sper$ is the symmetry parameter of the
  forbidden transition at the first prime $p = 3$ of the sieve
  (theorem~$\Tone$, \cite{PT_M1});
\item \textbf{geometric}: $\sper$ is the threshold $\gammap{p} > 1/2$
  for activity of a prime in the cascade
  (theorem~$\Tfive$, \cite{PT_M3});
\item \textbf{spectral}: $\sper$ is the real part appearing in the
  Mellin transform on the primes via the Jacobian of the multiplicative
  Haar measure.
\end{itemize}

A central reading of the present work is that these three manifestations
are operator-theoretically identifiable (\S\ref{sec:localisation_critique}).

\subsection{Fixed Point, Active Primes, Bifurcation}
\label{sec:point_fixe}

The PT sieve admits a unique fixed point $\mustar = 15$. The
arithmetically active primes at this fixed point are the prime divisors
$\{3, 5, 7\}$ of $\mustar$; the prime $p = 2$ plays the distinct role of
parity channel (cf.~\S\ref{sec:dualite_higgs}). At this fixed point, the
persistence dynamics bifurcates into two conjugate branches $\qplus$
and $\qminus$ characterised by
\begin{equation}
\qplus \;=\; \dfrac{13}{15}, \qquad
\qminus \;=\; \exp\!\Big(-\dfrac{1}{15}\Big).
\label{eq:qplusqminus}
\end{equation}
This bifurcation is the central object of the zeta programme. It enters
all the eigen-amplitudes and all the cohomological phases encountered
below.

\subsection{Multiplicative Haar Measure and Mellin Coordinate}
\label{sec:haar}

The persistence cascade singles out a single invariant measure on the
half-line of primes: the multiplicative Haar measure
\begin{equation}
du \;=\; \dfrac{dp}{p}, \qquad u \;=\; \log p.
\label{eq:haar_mult}
\end{equation}
This measure is invariant under dilation $p \mapsto \lambda p$, and it is
the unique measure invariant under the RG flow of PT (theorem~$\Ttwo$,
double stochasticity of the matrix $T_3$, \cite{PT_M1}). The
coordinate $u = \log p$ will be called the \emph{Mellin coordinate}
throughout the article; it carries the measure additively invariant under
translation $u \mapsto u + c$.

\subsection{Canonical PT Factorisation of $\zeta$}
\label{sec:factorisation}

For $\Reop(s) > 1$ we define, from the bifurcation amplitudes,
\begin{equation}
\delta_p^{\pm} \;=\; \dfrac{1 - q_{\pm}^{p}}{p}, \qquad
a_p \;=\; \delta_p^{+}\,(2 - \delta_p^{+}), \qquad
b_p \;=\; \delta_p^{-}\,(2 - \delta_p^{-}),
\label{eq:apbp}
\end{equation}
then
\begin{equation}
\zetaplus(s) \;=\; \exp\!\Big(\sum_{p} a_p\,p^{-s}\Big), \qquad
\zetaminus(s) \;=\; \exp\!\Big(\sum_{p} b_p\,p^{-s}\Big),
\label{eq:zetapm}
\end{equation}
and the \textbf{residual factor}
\begin{equation}
\Rres(s) \;=\; \dfrac{\zeta(s)}{\zetaplus(s)\,\zetaminus(s)}.
\label{eq:Rres_def}
\end{equation}
We finally set $\Ap = a_p + b_p$, which will be the effective PT
amplitude per prime in the sequel. Asymptotically, $\Ap \sim 4/p$ for
$p \to \infty$.

\paragraph{Canonical writing in holonomy angles.}
The amplitudes $a_p, b_p$ also write
\begin{equation}
a_p \;=\; \sinp{p}(\qplus), \qquad b_p \;=\; \sinp{p}(\qminus),
\label{eq:apbp_sinp}
\end{equation}
that is
\begin{equation}
\Ap \;=\; \sinp{p}(\qplus) \,+\, \sinp{p}(\qminus),
\label{eq:Ap_holonomie}
\end{equation}
where $\thetap{p}$ denotes the holonomy angle of the cascade at prime $p$
(monograph ch.~10, fine structure, \cite{PT_Foundations}). The equivalence
between the two writings comes from $1 - \delta_p^{\pm} = q_{\pm}^{p}/p$,
that is
\[
\delta_p^{\pm}\,(2 - \delta_p^{\pm}) \;=\; 1 - (1 - \delta_p^{\pm})^{2} \;=\; 1 - (q_{\pm}^{p}/p)^{2} \;=\; \sinp{p}(q_{\pm})
\]
on setting $\cos(\thetap{p}, q_{\pm}) = q_{\pm}^{p}/p$.

The zeta programme seeks a canonical spectral model of $\Rres$, under the
\textbf{conjectural hypothesis} of non-vanishing of
$\zetaplus\,\zetaminus$ (verified numerically on the tested zones) which
identifies the non-trivial zeros of $\Rres$ with those of $\zeta$. This
hypothesis will be recalled at each critical step of the identification.

